\documentclass[11pt,a4paper]{article}

\usepackage[T1]{fontenc}
\usepackage[utf8]{inputenc}
\usepackage[french]{babel}
\usepackage{lmodern}
\usepackage{microtype}
\DisableLigatures{encoding=T1,family=lmtt}
\usepackage[margin=23mm,headheight=15pt]{geometry}
\usepackage{amsmath,amssymb,mathtools}
\usepackage{stmaryrd}
\usepackage{xcolor}
\usepackage{array,booktabs,tabularx}
\usepackage{enumitem}
\usepackage{listings}
\usepackage[most]{tcolorbox}
\usepackage{fancyhdr}
\usepackage{titlesec}
\usepackage{hyperref}

\definecolor{ink}{HTML}{163044}
\definecolor{accent}{HTML}{087F8C}
\definecolor{pale}{HTML}{EFF7F8}
\definecolor{warm}{HTML}{FFF6E9}
\definecolor{muted}{HTML}{536675}
\definecolor{codebg}{HTML}{F4F6F8}

\hypersetup{colorlinks=true,linkcolor=accent,urlcolor=accent,
  pdftitle={CatOpt — Guide de démarrage},
  pdfauthor={Guide préparé pour Samir},
  pdfsubject={Apprendre le C++20 par les morphismes, la composition et la fonctorialité}}
\setlength{\parindent}{0pt}
\setlength{\parskip}{6pt}
\clubpenalty=10000
\widowpenalty=10000
\setlist[enumerate]{leftmargin=*,itemsep=5pt,topsep=5pt}
\setlist[itemize]{leftmargin=*,itemsep=3pt,topsep=4pt}
\setcounter{tocdepth}{1}
\setcounter{secnumdepth}{0}
\titleformat{\section}{\Large\bfseries\color{ink}}{}{0pt}{}
\titleformat{\subsection}{\large\bfseries\color{ink}}{}{0pt}{}
\titlespacing*{\section}{0pt}{8pt}{9pt}
\titlespacing*{\subsection}{0pt}{12pt}{4pt}
\pagestyle{fancy}
\fancyhf{}
\fancyhead[L]{\small\color{muted}\textbf{CatOpt} / Apprendre le C++ par les catégories}
\fancyfoot[L]{\small\color{muted}Guide de démarrage · C++20}
\fancyfoot[R]{\small\thepage}
\renewcommand{\headrulewidth}{0.3pt}
\renewcommand{\footrulewidth}{0pt}
\newcommand{\code}[1]{\texttt{\detokenize{#1}}}
\newcommand{\Id}{\operatorname{id}}
\newcommand{\List}{\operatorname{List}}
\newcommand{\sem}[1]{\llbracket #1\rrbracket}
\newcommand{\checkpoint}{\(\square\)\ }
\newcommand{\etape}[2]{\clearpage\section{Étape #1 — #2}}

\tcbset{enhanced,boxrule=0.4pt,arc=2mm,
  left=3mm,right=3mm,top=2mm,bottom=2mm,
  before skip=8pt,after skip=8pt,fonttitle=\bfseries}
\newtcolorbox{objectif}{colback=pale,colframe=accent,
  coltitle=ink,colbacktitle=pale,title={Pourquoi cette étape existe dans CatOpt}}
\newtcolorbox{theorie}{colback=white,colframe=accent,
  coltitle=white,colbacktitle=accent,title={Le lien avec la théorie des catégories}}
\newtcolorbox{validation}{colback=codebg,colframe=muted,
  coltitle=ink,colbacktitle=codebg,title={Tu peux passer à la suite lorsque…}}
\newtcolorbox{vigilance}{colback=warm,colframe=orange!55!black,
  coltitle=ink,colbacktitle=warm,title={Point à comprendre}}

\lstdefinestyle{cpp}{language=C++,basicstyle=\small\ttfamily,
  keywordstyle=\color{accent}\bfseries,commentstyle=\color{muted},
  stringstyle=\color{orange!65!black},backgroundcolor=\color{codebg},
  frame=single,rulecolor=\color{codebg},framerule=0.3pt,
  columns=fullflexible,keepspaces=true,showstringspaces=false,
  breaklines=true,tabsize=4,xleftmargin=2mm,xrightmargin=2mm,
  aboveskip=8pt,belowskip=8pt,upquote=true,
  morekeywords={uint64_t,size_t}}
\lstset{style=cpp}
\lstdefinestyle{plain}{style=cpp,language={},keywordstyle={},
  commentstyle={},stringstyle={}}
\BeforeBeginEnvironment{lstlisting}{\par\noindent\begin{minipage}{\linewidth}}
\AfterEndEnvironment{lstlisting}{\end{minipage}\par}

\begin{document}

\hypersetup{pageanchor=false}
\begin{titlepage}
\vspace*{13mm}
{\color{accent}\large CARNET DE CONSTRUCTION}\par
\vspace{9mm}
{\color{ink}\fontsize{46}{50}\selectfont\bfseries CatOpt}\par
\vspace{5mm}
{\color{ink}\LARGE Apprendre le C++\par avec la théorie des catégories}\par
\vspace{8mm}
{\large Un guide pas à pas pour construire ton premier\par langage de morphismes et son optimiseur.}
\vspace{13mm}

\begin{objectif}
Tu vas représenter des programmes comme des morphismes, les composer,
les exécuter, puis les transformer à l'aide d'équations explicites.
Chaque nouvelle notion de C++ répondra à un besoin de cette construction.
\end{objectif}

\vspace{6mm}
\[
 U \xrightarrow{\;A_1\;} U \xrightarrow{\;S\;} U
 \qquad\leadsto\qquad
 \List(U)\xrightarrow{\;F(S\circ A_1)\;}\List(U)
\]
\begin{center}
\color{muted}Ajouter un, élever au carré, puis transformer une liste\par
avec un seul parcours au lieu de deux.
\end{center}

\vfill
\textbf{Point de départ :} tu as déjà pratiqué les classes et rencontré les
smart pointers ; les bases restent à consolider et les templates à découvrir.

\textbf{Format :} tâches courtes, signatures proposées, exemples attendus,
questions de compréhension et critères de réussite. Les fonctions centrales
restent à écrire par toi.

\vspace{5mm}
{\small\color{muted}Volume 1 · Du premier morphisme à la fusion des \code{map}\par
Version du 6 novembre 2026}
\end{titlepage}

\pagenumbering{roman}
\hypersetup{pageanchor=true}
\tableofcontents
\clearpage
\pagenumbering{arabic}
\section{Le fil conducteur : que construit-on exactement ?}

L'objectif principal est de progresser en C++ en construisant quelque chose
qui te plaît. CatOpt sera d'abord un \textbf{petit langage de calculs purs}
écrit en C++. Il manipulera des descriptions de programmes que tu peux
inspecter. Ce volume aboutit à un interpréteur et à deux petites optimisations ;
la génération de code C++ viendra dans un volume ultérieur.

Le chemin complet sera :
\[
\boxed{\text{Décrire}}\;\longrightarrow\;
\boxed{\text{Composer}}\;\longrightarrow\;
\boxed{\text{Interpréter}}\;\longrightarrow\;
\boxed{\text{Réécrire}}.
\]

\subsection{Le résultat concret que tu vises}
Tu construiras le morphisme \(f(x)=x+1\), le morphisme \(g(x)=x^2\), puis
leur composition \(g\circ f\). Tu pourras afficher sa description et
vérifier que l'entrée \(3\) produit \(16\).

Ensuite, tu compareras ces deux programmes de listes :
\[
\underbrace{F(f)\text{ puis }F(g)}_{\text{deux parcours}}
\qquad\text{et}\qquad
\underbrace{F(g\circ f)}_{\text{un parcours}}.
\]
Ils donneront les mêmes valeurs ; le second évitera un tableau intermédiaire
dans l'interpréteur que tu auras construit.

\subsection{Pourquoi commencer par des suites d'opérations ?}
Un morphisme sera d'abord une suite comme
\code{[Add(1), Square]}. Tu apprendras ainsi les valeurs, les références,
les conteneurs et la composition sans devoir gérer immédiatement un arbre
récursif. Cette représentation a une raison mathématique : un morphisme
librement construit à partir de générateurs peut être représenté par un chemin.

Nous commencerons avec \textbf{un seul objet} \(U\), donc avec des morphismes
\(U\to U\), appelés \emph{endomorphismes}. C'est déjà une catégorie.
Le fait que toutes les compositions soient typées correctement découlera
de cette restriction. Le vérificateur de types général arrivera quand nous
introduirons d'autres domaines et codomaines.

La notation \(\sem{p}(x)\) signifiera « le résultat de l'interprétation
du programme \(p\) sur la valeur \(x\) ». Elle distingue la description
\(p\) de la fonction qu'elle représente.

\subsection{Comment utiliser ce guide}
Pour chaque étape : lis l'objectif, écris sur papier l'exemple mathématique,
implémente une seule tâche, compile, puis exécute les vérifications.
Si tu bloques, isole une question précise : « Que contient ce vecteur ? »,
« Qui possède cette valeur ? », « Dans quel ordre les fonctions s'appliquent-elles ? »

\begin{vigilance}
Les blocs de signatures indiquent les fonctions à écrire : ce sont des
\emph{déclarations}, pas des implémentations complètes. Une déclaration qui
finit par un point-virgule indique au compilateur qu'une fonction existe.
Sa définition, avec un corps entre accolades, sera ton travail.
\end{vigilance}

Garde une courte note après chaque étape : ce que tu as construit, la loi
mathématique concernée, une notion C++ comprise et une question encore ouverte.
Tu peux répartir une étape sur plusieurs séances. Le critère de progression
est ta compréhension, pas la vitesse.

\etape{0}{Préparer un atelier minimal}
\label{sec:atelier}
\begin{objectif}
Pouvoir compiler et exécuter une petite modification immédiatement.
Ce cycle court te permettra de confronter une définition mathématique à son
implémentation sans perdre le fil du projet.
\end{objectif}

\textbf{C++ à travailler :} point d'entrée \code{main}, en-têtes, compilation,
avertissements et assertions.

\subsection{À implémenter, dans cet ordre}
\begin{enumerate}
\item Dans \code{main.cpp}, écris un programme qui affiche
\code{CatOpt - premier morphisme}. Il doit inclure \code{<iostream>}.
\item Crée un dossier \code{build}, puis compile depuis la racine du projet.
Avec Clang, les commandes sont :
\begin{lstlisting}[style=plain]
mkdir -p build
clang++ -std=c++20 -Wall -Wextra -Wpedantic -g -O0 \
    main.cpp -o build/catopt
./build/catopt
\end{lstlisting}
Avec GCC, remplace \code{clang++} par \code{g++}. Le dossier \code{build}
contient les exécutables ; ton code source reste à la racine pour l'instant.
\item Ajoute \code{<cassert>} et une assertion vraie dans \code{main}.
Change-la volontairement en assertion fausse, observe l'arrêt du programme,
puis rétablis-la. Tu sais maintenant reconnaître un test qui échoue.
\item Crée une fonction \code{void run_tests();}, appelle-la depuis
\code{main}, et place tes assertions dans sa définition. Elle accueillera
les premiers tests du langage.
\end{enumerate}

\begin{theorie}
Nous voulons vérifier des propriétés des morphismes. L'atelier de tests sera
le lieu où tu écriras des exemples de ces propriétés. Une assertion n'est pas
une preuve universelle : elle vérifie un cas concret et révèle des erreurs
dans l'implémentation.
\end{theorie}

\subsection{Questions à te poser}
Que fait \code{#include} ? Quelle différence fais-tu entre une erreur de
compilation, une erreur d'édition de liens et un échec à l'exécution ?
Peux-tu retrouver le premier message utile quand une compilation échoue ?

\begin{validation}
\checkpoint Tu sais recompiler après une modification.\par
\checkpoint Le programme s'exécute et une assertion fausse l'arrête.\par
\checkpoint Tu as compris chaque avertissement restant, et corrigé ceux qui
signalent une erreur.\par
Pour cette phase, conserve les assertions actives : ne définis pas
\code{NDEBUG}, qui les désactive.
\end{validation}

\etape{1}{Donner un sens précis aux calculs}
\label{sec:scalaires}
\begin{objectif}
Définir les valeurs que CatOpt manipule et les premières fonctions qu'il
saura représenter. Une optimisation ne peut être justifiée que si l'on sait
précisément ce que signifie le programme initial.
\end{objectif}

\textbf{C++ à travailler :} alias de type, paramètres, retour par valeur,
arithmétique non signée et limites d'un type.

\begin{theorie}
Notre premier objet sera \(U=\mathbb{Z}/2^{64}\mathbb{Z}\), représenté par les
entiers de \(0\) à \(2^{64}-1\). Les calculs se font modulo \(2^{64}\).
Nous définissons trois familles de fonctions :
\[
\Id_U(x)=x,\qquad A_c(x)=x+c\pmod{2^{64}},\qquad
S(x)=x^2\pmod{2^{64}}.
\]
Toutes ont le type mathématique \(U\to U\).
\end{theorie}

\subsection{À implémenter, dans cet ordre}
\begin{enumerate}
\item Inclus \code{<cstdint>} et déclare :
\begin{lstlisting}
using Scalar = std::uint64_t;
\end{lstlisting}
Cet alias donne un nom de projet au type. Il ne crée pas un nouveau type
distinct. Ce guide suppose une plateforme qui fournit \code{std::uint64_t},
comme les plateformes usuelles de Clang et GCC.
\item Écris les trois fonctions suivantes :
\begin{lstlisting}
Scalar identity_value(Scalar x);
Scalar add_value(Scalar x, Scalar constant);
Scalar square_value(Scalar x);
\end{lstlisting}
\item Dans les tests, compose manuellement l'addition de \(1\) et le carré.
L'entrée \(3\) doit donner \(16\). Inverse l'ordre : tu dois obtenir \(10\).
\item Inclus \code{<limits>}. Obtiens la valeur maximale avec
\code{std::numeric_limits<Scalar>::max()} et vérifie que lui ajouter
\code{Scalar{1}} produit zéro.
\end{enumerate}

\subsection{Exemples attendus}
\begin{center}
\begin{tabular}{lll}
\toprule
Calcul & Entrée & Résultat \\
\midrule
\(\Id_U\) & \(7\) & \(7\) \\
\(A_1\) & \(3\) & \(4\) \\
\(S\) & \(4\) & \(16\) \\
\(S\circ A_1\) & \(3\) & \(16\) \\
\(A_1\circ S\) & \(3\) & \(10\) \\
\(A_1\) & \(2^{64}-1\) & \(0\) \\
\bottomrule
\end{tabular}
\end{center}

\begin{vigilance}
Le choix du type fait partie de la sémantique. Ici, addition et multiplication
sur nos valeurs non signées de 64 bits suivent l'arithmétique modulaire.
Nous ne prétendons pas calculer sur tous les entiers mathématiques.
Employer un \code{int} signé et laisser le résultat dépasser son domaine
introduirait un comportement indéfini en C++.
\end{vigilance}

Les fonctions du langage sont pures : pas d'affichage, de mutation d'une
variable globale ni d'aléatoire dans leur calcul. Les diagnostics de CatOpt
seront réalisés autour de l'évaluation. Nous raisonnons sur les résultats
mathématiques ; les échecs matériels ou d'allocation sont en dehors du modèle
de ce premier volume.

\subsection{Questions à te poser}
Pourquoi passer un petit entier par valeur ? Quelle différence y a-t-il entre
le paramètre \code{x}, une variable locale et la valeur retournée ?
Pourquoi l'ordre de composition change-t-il le résultat ?

\begin{validation}
\checkpoint Tous les exemples du tableau sont vérifiés.\par
\checkpoint Tu peux expliquer ce qu'est l'objet \(U\).\par
\checkpoint Tu peux expliquer pourquoi \(A_1\circ S\neq S\circ A_1\).
\end{validation}

\textbf{Lien avec la suite :} pour l'instant, tu exécutes directement des
fonctions C++. L'optimiseur aura besoin d'en manipuler la \emph{description}.

\etape{2}{Représenter les générateurs comme des données}
\label{sec:primitives}
\begin{objectif}
Permettre à CatOpt d'inspecter une opération avant de l'exécuter. Pour
reconnaître « ajouter zéro », le moteur doit voir le nom de l'opération et
sa constante, pas seulement disposer d'une fonction exécutable.
\end{objectif}

\textbf{C++ à travailler :} \code{enum class}, \code{struct}, initialisation,
\code{switch}, \code{const&}, \code{std::string}.

\begin{theorie}
Les opérations \(A_c\) et \(S\) sont nos \emph{générateurs}. Un générateur est
un bloc élémentaire à partir duquel nous construirons des morphismes par
composition. Sa description est syntaxique ; la fonction \(U\to U\) qu'il
désigne est son interprétation.
\end{theorie}

\subsection{À implémenter, dans cet ordre}
\begin{enumerate}
\item Déclare un \code{enum class PrimitiveKind} contenant deux valeurs :
\code{AddConstant} et \code{Square}.
\item Déclare une \code{struct Primitive} avec un champ \code{kind} de ce type
et un champ \code{constant} de type \code{Scalar}.
\item Écris deux fonctions de construction :
\begin{lstlisting}
Primitive make_add(Scalar constant);
Primitive make_square();
\end{lstlisting}
Pour \code{Square}, fixe \code{constant} à zéro. Ce champ n'y sert pas, mais
une convention stable facilitera les comparaisons de descriptions.
\item Écris l'interpréteur d'une primitive :
\begin{lstlisting}
Scalar evaluate(const Primitive& primitive, Scalar x);
\end{lstlisting}
Utilise un \code{switch} sur \code{kind} et les fonctions de l'étape 1.
N'ajoute pas de retour arbitraire comme zéro pour masquer un cas invalide ;
tu peux terminer par une exception \code{std::logic_error} de
\code{<stdexcept>} si aucun cas reconnu n'a été rencontré.
\item Écris une fonction d'affichage qui retourne une chaîne :
\begin{lstlisting}
std::string to_string(const Primitive& primitive);
\end{lstlisting}
Elle doit produire \code{Add(3)} ou \code{Square}. Inclus \code{<string>}.
\end{enumerate}

\subsection{Ce que tu dois observer}
\begin{lstlisting}
const auto f = make_add(1);
const auto g = make_square();
assert(evaluate(f, 3) == 4);
assert(evaluate(g, 4) == 16);
assert(to_string(f) == "Add(1)");
\end{lstlisting}
\code{auto} laisse le compilateur déduire ici le type \code{Primitive}.
Le mot \code{const} interdit de modifier la variable après son initialisation.

\subsection{Pourquoi une référence constante ?}
Dans \code{const Primitive&}, le symbole \code{&} désigne une référence à
l'objet fourni par l'appelant. La fonction peut le lire sans en faire une
copie, et \code{const} lui interdit de le modifier par cette référence.
La primitive est petite : un passage par valeur serait aussi défendable.
Le choix prépare surtout le même contrat de lecture pour les morphismes,
qui contiendront plusieurs opérations.

\begin{vigilance}
La structure publique permet encore de construire des descriptions qui ne
respectent pas nos conventions. Pour ce premier noyau, construis les
primitives uniquement avec \code{make_add} et \code{make_square}.
À l'étape 8, tu rendras cette convention explicite dans l'interface.
\end{vigilance}

\subsection{Questions à te poser}
Quelle différence y a-t-il entre \code{make_add(1)} et
\code{add_value(3, 1)} ? Une description doit-elle contenir son entrée ?
Pourquoi l'affichage ne doit-il pas exécuter la primitive ?

\begin{validation}
\checkpoint Chaque primitive peut être construite, affichée et évaluée.\par
\checkpoint \code{evaluate} ne modifie pas la description reçue.\par
\checkpoint Tu distingues « construire le morphisme » et « calculer sa valeur ».
\end{validation}

\etape{3}{Construire un morphisme comme un chemin}
\label{sec:morphismes}
\begin{objectif}
Représenter un programme entier avec plusieurs générateurs. Cette
représentation sera l'entrée de l'interpréteur et, plus tard, de l'optimiseur.
\end{objectif}

\textbf{C++ à travailler :} \code{std::vector}, boucle sur une collection,
propriété par valeur, copie, surcharges et cas vide.

\begin{theorie}
Un chemin de générateurs donne un morphisme :
\[
U\xrightarrow{A_1}U\xrightarrow{S}U.
\]
Nous stockons les opérations dans \textbf{l'ordre d'exécution}, de gauche à
droite : \code{[Add(1), Square]} représente \(S\circ A_1\).
Le chemin vide représente \(\Id_U\), puisqu'il laisse son entrée inchangée.
\end{theorie}

\subsection{À implémenter, dans cet ordre}
\begin{enumerate}
\item Inclus \code{<vector>} et introduis la représentation :
\begin{lstlisting}
struct Morphism {
    std::vector<Primitive> steps;
};
\end{lstlisting}
\item Écris deux fonctions de construction :
\begin{lstlisting}
Morphism identity();
Morphism lift(Primitive primitive);
\end{lstlisting}
La première crée un chemin vide ; la seconde un chemin contenant exactement
la primitive fournie. Ici, \code{lift} signifie seulement « mettre une
primitive dans une description de morphisme » ; nous n'introduisons pas
une nouvelle structure catégorique avec ce nom.
\item Écris l'évaluation du chemin :
\begin{lstlisting}
Scalar evaluate(const Morphism& morphism, Scalar x);
\end{lstlisting}
Pars d'une valeur courante égale à \code{x}. Pour chaque primitive, remplace
la valeur courante par le résultat de son évaluation. Retourne la valeur
finale. Le chemin vide doit fonctionner sans traitement artificiel.
Les fonctions \code{evaluate} pour \code{Primitive} et pour \code{Morphism}
portent le même nom : ce sont des \emph{surcharges}. Le compilateur choisit
la fonction selon le type des arguments. Le même principe s'applique à
\code{to_string}.
\item Écris \code{std::string to_string(const Morphism& morphism);}.
Affiche \code{Id} pour le chemin vide et, par exemple,
\code{Add(1) >> Square} pour les deux étapes précédentes.
\item Copie un morphisme dans une deuxième variable, modifie les étapes de
la copie et vérifie que l'original n'a pas changé.
\end{enumerate}

\subsection{Exemples attendus}
\begin{lstlisting}
const Morphism p{{make_add(1), make_square()}};
assert(evaluate(p, 3) == 16);
assert(evaluate(identity(), 3) == 3);
assert(evaluate(lift(make_square()), 5) == 25);
\end{lstlisting}
Les deux niveaux d'accolades initialisent la structure puis son vecteur.
Si cette syntaxe te gêne, commence avec une variable vide et deux appels
à \code{push_back} ; vérifie ensuite que les deux écritures décrivent le même chemin.

\subsection{Pourquoi aucun smart pointer pour le moment ?}
Le \code{Morphism} possède son vecteur, qui possède ses primitives. Quand le
morphisme est détruit, les ressources de son vecteur sont libérées
automatiquement. C'est déjà une application du principe RAII : la durée de
vie des ressources est attachée à celle des objets.

Un \code{std::shared_ptr} introduirait ici une question de partage que notre
modèle n'exige pas. Les pointeurs intelligents deviendront pertinents si
tu choisis ensuite de représenter des expressions récursives ou partagées.

\subsection{Questions à te poser}
Quel objet possède les primitives ? Que copie-t-on en copiant un
\code{Morphism} ? Pourquoi la variable courante de \code{evaluate} peut-elle
changer alors que le morphisme reçu est constant ?

\begin{validation}
\checkpoint L'identité est un vecteur vide et conserve toute entrée testée.\par
\checkpoint Un chemin est évalué dans son ordre de stockage.\par
\checkpoint Modifier une copie ne modifie pas l'original.\par
\checkpoint Tu peux relier le vecteur au chemin de morphismes dessiné plus haut.
\end{validation}

\etape{4}{Composer et vérifier les lois de catégorie}
\label{sec:composition}
\begin{objectif}
Disposer de l'opération centrale du langage : construire un programme à
partir de deux programmes. La composition servira ensuite à fusionner les
calculs effectués sur les éléments d'une liste.
\end{objectif}

\textbf{C++ à travailler :} lecture par référence constante, création d'un
résultat par valeur, concaténation et égalité de descriptions.

\begin{theorie}
La convention est fixe : \code{compose(g, f)} représente \(g\circ f\), donc
\textbf{on exécute d'abord \(f\), puis \(g\)}. Les lois attendues sont :
\[
\Id_U\circ f=f,\qquad f\circ\Id_U=f,\qquad
h\circ(g\circ f)=(h\circ g)\circ f.
\]
Elles portent sur la composition. L'associativité de la composition ne
signifie pas que toutes les opérations arithmétiques sont associatives.
\end{theorie}

\subsection{À implémenter, dans cet ordre}
\begin{enumerate}
\item Écris :
\begin{lstlisting}
Morphism compose(const Morphism& g, const Morphism& f);
\end{lstlisting}
Construis un résultat contenant les étapes de \code{f}, puis celles de
\code{g}. Commence par deux boucles avec \code{push_back}. Tu pourras
ensuite étudier \code{insert} et \code{reserve}, sans changer le contrat.
\item Vérifie que la composition ne modifie aucun argument.
Retourne le résultat par valeur. Ne retourne jamais une référence vers une
variable locale, dont la durée de vie se termine à la sortie de la fonction.
\item Écris une comparaison structurelle explicite :
\begin{lstlisting}
bool same_syntax(const Morphism& a, const Morphism& b);
\end{lstlisting}
Compare la longueur, les sortes d'opérations et, pour une addition, la
constante. Cette fonction compare des descriptions, pas toutes les valeurs
qu'elles pourraient produire.
\item Vérifie les deux lois d'identité et l'associativité avec cette
comparaison, sur plusieurs chemins, dont le chemin vide.
\item Vérifie aussi l'accord entre composition et interprétation :
\begin{lstlisting}
const auto f = lift(make_add(1));
const auto g = lift(make_square());
const Scalar x{3};
assert(evaluate(compose(g, f), x)
       == evaluate(g, evaluate(f, x)));
\end{lstlisting}
Utilise notamment \code{x} égal à \(0\), \(3\), \(7\) et à la valeur maximale.
\end{enumerate}

\subsection{Une erreur d'ordre que tes tests doivent détecter}
Pose \(f=A_1\) et \(g=S\). Tu dois obtenir :
\[
\sem{\operatorname{compose}(g,f)}(3)=16,
\qquad
\sem{\operatorname{compose}(f,g)}(3)=10.
\]
Tester seulement deux additions ne suffirait pas à détecter une inversion :
elles commutent ici. Choisir un exemple discriminant fait partie de la
conception d'un bon test.

\subsection{Pourquoi les lois sont-elles déjà vraies dans la structure ?}
La composition concatène des suites. Ajouter la suite vide ne change rien ;
concaténer trois suites donne les mêmes éléments dans le même ordre, quel
que soit le parenthésage. Notre représentation a déjà oublié les parenthèses
inutiles des compositions.

Les chemins forment la catégorie libre à un objet engendrée par nos
primitives. « Libre » signifie ici que l'on n'a pas encore identifié les
chemins au moyen de relations supplémentaires entre générateurs.

\begin{vigilance}
\code{[Add(1), Add(2)]} et \code{[Add(3)]} sont deux chemins différents,
mais ils désignent la même fonction sur \(U\). L'interprétation
\(p\mapsto\sem{p}\) n'est donc pas injective. L'optimiseur exploitera
précisément certaines de ces égalités sémantiques.
\end{vigilance}

\begin{validation}
\checkpoint Tu sais expliquer l'ordre des arguments de \code{compose}.\par
\checkpoint Les trois lois sont vérifiées structurellement sur tes exemples.\par
\checkpoint Tu peux justifier ces lois par la concaténation, au-delà des tests.\par
\checkpoint Tu distingues égalité de syntaxe et égalité de fonctions.
\end{validation}

\etape{5}{Écrire les premières réécritures justifiées}
\label{sec:reecriture}
\begin{objectif}
Transformer une description en une autre, plus courte, sans changer la
fonction calculée. CatOpt devient ici un premier optimiseur.
\end{objectif}

\textbf{C++ à travailler :} parcours et construction d'un vecteur, cas limites,
invariant de boucle et contrat d'une fonction.

\begin{theorie}
Nous ajoutons deux familles d'équations :
\[
A_0=\Id_U,\qquad A_b\circ A_a=A_{a+b},
\]
où \(a+b\) est calculé modulo \(2^{64}\). Ce sont des lois de nos
\emph{générateurs arithmétiques}, pas des lois vraies pour tous les morphismes
d'une catégorie quelconque.

Si \(E\) désigne ces équations, les chemins peuvent être regroupés modulo
les égalités qu'elles engendrent, y compris à l'intérieur d'une composition.
On obtient un quotient de la catégorie syntaxique. Pour construire ce
premier moteur, tu n'as pas besoin d'implémenter les classes du quotient.
\end{theorie}

\subsection{Écris d'abord la justification sur papier}
Pour toute valeur \(x\in U\),
\[
\sem{A_b\circ A_a}(x)=(x+a)+b=x+(a+b)=\sem{A_{a+b}}(x).
\]
L'associativité de l'addition modulaire justifie cette transformation.
Pour \(A_0\), la justification est \(x+0=x\).

Ces égalités restent valides à l'intérieur d'un chemin : remplacer une
sous-composition par une fonction égale préserve la composition avec ce
qui vient avant et après. C'est la compatibilité de l'égalité avec la composition.

\subsection{À implémenter, dans cet ordre}
\begin{enumerate}
\item Déclare :
\begin{lstlisting}
Morphism simplify(const Morphism& input);
\end{lstlisting}
La fonction retourne une nouvelle description ; elle ne modifie pas
\code{input}. Première version : reconstruis simplement le chemin en
omettant chaque \code{Add(0)}.
\item Ajoute la fusion d'additions \emph{adjacentes}. Une stratégie possible
est d'accumuler une constante en attente pendant que tu lis une série
d'additions. Quand tu rencontres \code{Square}, émets l'addition accumulée
si elle est non nulle, remets l'accumulateur à zéro, puis émets \code{Square}.
\item N'oublie pas d'émettre l'addition encore en attente à la fin du chemin.
Si sa constante est zéro, aucune primitive ne doit être ajoutée.
\item Vérifie que \code{simplify} ne modifie pas son entrée et préserve les
valeurs calculées sur les exemples ci-dessous.
\item Vérifie qu'une deuxième simplification produit la même syntaxe que
la première : c'est l'idempotence de \code{simplify}.
\end{enumerate}

\subsection{Exemples attendus}
Dans ce tableau, \(M=2^{64}-1\) et \code{[]} désigne le chemin vide.
\begin{center}
\small
\begin{tabularx}{\linewidth}{@{}X X@{}}
\toprule
Avant & Après \\
\midrule
\code{[]} & \code{[]} \\
\code{[Add(0)]} & \code{[]} \\
\code{[Add(1), Add(2)]} & \code{[Add(3)]} \\
\code{[Add(1), Add(2), Add(3)]} & \code{[Add(6)]} \\
\code{[Add(1), Square, Add(2)]} & Inchangé \\
\code{[Add(M), Add(1)]} & \code{[]} \\
\code{[Add(0), Square, Add(2), Add(3)]} & \code{[Square, Add(5)]} \\
\bottomrule
\end{tabularx}
\end{center}

Sur l'entrée \(3\), le dernier chemin vaut \(14\) avant comme après
simplification. Vérifie aussi plusieurs entrées proches de la borne du type.

\subsection{L'invariant qui aide à programmer}
Après avoir lu une partie du chemin, la sortie déjà construite, suivie de
l'addition encore en attente, doit calculer la même fonction que la partie
lue. Cet énoncé explique pourquoi tu peux différer l'émission des additions.
Il t'aide aussi à repérer deux erreurs classiques : oublier la fin du chemin
ou laisser une addition traverser un carré.

\begin{vigilance}
Tu ne peux pas déplacer librement une addition à travers \code{Square}.
Sur \(3\), ajouter \(1\) puis élever au carré donne \(16\), tandis que
l'ordre inverse donne \(10\). Une transformation doit avoir sa propre
justification ; la présence d'une composition ne permet pas de permuter ses termes.
\end{vigilance}

\subsection{Mesurer ce que cette optimisation fait réellement}
Définis le coût syntaxique \(C(p)\) comme le nombre de primitives du chemin
\(p\). Vérifie sur tes exemples que :
\[
C(\operatorname{simplify}(p))\leq C(p).
\]
Cela décrit une réduction d'opérations interprétées. Cela ne démontre
encore aucun gain de temps sur un programme compilé.

Cette fonction n'est pas un décideur de toutes les égalités entre fonctions
sur \(U\). Elle applique seulement les règles que tu lui as données.

\begin{validation}
\checkpoint Les exemples, notamment le chemin vide et le retour modulaire à
zéro, sont traités correctement.\par
\checkpoint Tu justifies les deux règles pour toute entrée \(x\).\par
\checkpoint Tes tests comparent avant/après et vérifient l'idempotence.\par
\checkpoint Tu sais distinguer cette justification générale des tests de ton code.
\end{validation}

\textbf{Premier jalon atteint :} tu disposes d'un langage, d'une interprétation
et d'une transformation qui préserve cette interprétation.

\etape{6}{Passer des scalaires aux listes avec un foncteur}
\label{sec:map}
\begin{objectif}
Réutiliser chaque morphisme scalaire pour transformer tous les éléments
d'une liste. Cette construction prépare directement la fusion des parcours.
\end{objectif}

\textbf{C++ à travailler :} vecteur d'éléments, entrée en lecture seule,
construction d'une sortie, \code{reserve} et fonctions à plusieurs surcharges.

\begin{theorie}
Le foncteur des listes finies associe à un ensemble \(A\) l'ensemble
\(\List(A)\), et à une fonction \(f:A\to B\) la fonction :
\[
F(f)([x_1,\ldots,x_n])=[f(x_1),\ldots,f(x_n)].
\]
Les deux lois de foncteur sont :
\[
F(\Id_U)=\Id_{\List(U)},\qquad
F(g\circ f)=F(g)\circ F(f).
\]
\end{theorie}

\subsection{À implémenter, dans cet ordre}
\begin{enumerate}
\item Déclare un alias \code{using Values = std::vector<Scalar>; }.
Il représente les listes de valeurs de notre langage.
\item Écris :
\begin{lstlisting}
Values map_values(const Morphism& f, const Values& input);
\end{lstlisting}
Crée une sortie vide, lis chaque élément de l'entrée, évalue \code{f} sur
cet élément et ajoute le résultat à la sortie.
\item Vérifie que la liste vide est acceptée et que l'entrée n'est jamais
modifiée. L'ordre et le nombre des éléments doivent être conservés.
\item Ajoute \code{reserve(input.size())} sur le vecteur de sortie pour
réserver sa capacité. Cela ne crée aucun élément : ne confonds pas
\code{reserve} avec \code{resize}.
\item Écris des tests pour les deux lois de foncteur. Tu peux comparer
directement les vecteurs de scalaires avec \code{==}.
\end{enumerate}

\subsection{Exemples attendus}
Avec \(f=A_1\), \(g=S\) et \(xs=[0,1,2,3]\),
\[
F(f)(xs)=[1,2,3,4],\qquad
F(g)(F(f)(xs))=[1,4,9,16].
\]
Les assertions centrales auront cette forme :
\begin{lstlisting}
const Values xs{0, 1, 2, 3};
assert(map_values(identity(), xs) == xs);
assert(map_values(compose(g, f), xs)
       == map_values(g, map_values(f, xs)));
\end{lstlisting}
Répète ces vérifications avec une liste vide, un singleton et une liste
contenant la valeur maximale du type.

\subsection{La justification tient élément par élément}
Pour un indice \(i\), le côté gauche de la seconde loi produit
\(\sem{g\circ f}(x_i)=\sem{g}(\sem{f}(x_i))\). Le côté droit produit
exactement la même valeur à la même position. Les listes ont la même
longueur ; elles sont donc égales. La liste vide satisfait également la loi.

\begin{vigilance}
La loi \(F(\Id_U)=\Id_{\List(U)}\) exprime une égalité de valeurs produites.
Appliquer \code{map_values} à l'identité peut construire une nouvelle liste
égale à l'entrée. La loi n'exige ni la même adresse mémoire ni l'absence
d'allocation. Ce décalage entre résultat et coût est justement ce que
l'optimiseur pourra exploiter.
\end{vigilance}

\subsection{Précision sur les catégories employées}
On peut voir cette construction comme un foncteur de la catégorie des
chemins scalaires vers une catégorie à un objet \(\List(U)\), dont les
morphismes sont les fonctions de listes ainsi obtenues. L'égalité dans la
cible est l'égalité des fonctions, pas celle de leurs descriptions.

C'est aussi l'interprétation scalaire suivie du foncteur des listes
\(\List:\mathbf{Set}\to\mathbf{Set}\). Nous ne prétendons pas avoir un
endofoncteur sur une petite catégorie possédant seulement \(U\) et
\(\List(U)\) : appliquer de nouveau la construction ferait intervenir
\(\List(\List(U))\), que notre langage ne représente pas encore.

\begin{validation}
\checkpoint Les valeurs, l'ordre, la longueur et le cas vide sont corrects.\par
\checkpoint Tu sais expliquer ce que \(F\) fait à l'objet et aux morphismes.\par
\checkpoint Tu peux démontrer les lois élément par élément.\par
\checkpoint Tu distingues égalité du résultat et coût de l'exécution.
\end{validation}

\etape{7}{Fusionner les parcours par fonctorialité}
\label{sec:fusion}
\begin{objectif}
Réaliser la première optimisation directement motivée par une loi de
foncteur : remplacer plusieurs parcours d'une liste par un seul parcours
qui exécute le morphisme composé sur chaque élément.
\end{objectif}

\textbf{C++ à travailler :} composition de structures, deux niveaux de
parcours, séparation entre transformation d'un programme et exécution.

\begin{theorie}
La loi employée est exactement :
\[
\boxed{F(g)\circ F(f)=F(g\circ f).}
\]
Les deux côtés désignent la même fonction sur les listes. Ils ont cependant
deux descriptions et deux stratégies d'exécution différentes dans CatOpt.
Le moteur choisira celle qui matérialise moins de listes intermédiaires.
\end{theorie}

\subsection{7.1 — Décrire plusieurs parcours successifs}
Introduis une structure distincte du morphisme scalaire :
\begin{lstlisting}
struct MapPipeline {
    std::vector<Morphism> stages;
};
\end{lstlisting}
Chaque élément du vecteur représente \emph{un parcours complet} de la liste.
Le pipeline \code{[f, g]} signifie : appliquer \code{map_values(f, ...)},
puis \code{map_values(g, ...)}. Les étapes sont stockées dans l'ordre
d'exécution, comme les primitives d'un morphisme.

Écris les fonctions suivantes :
\begin{lstlisting}
Values evaluate(const MapPipeline& pipeline, const Values& input);
std::string to_string(const MapPipeline& pipeline);
\end{lstlisting}
Pour l'évaluation, commence avec une copie de l'entrée, puis remplace la
liste courante par le résultat de chaque \code{map_values}.
Un pipeline vide retourne donc une liste égale à l'entrée.
Affiche \code{IdList} pour le pipeline vide et, par exemple,
\code{Map(Add(1)) >> Map(Square)} pour deux parcours.

\newpage
\subsection{7.2 — Construire la description fusionnée}
Déclare :
\begin{lstlisting}
MapPipeline fuse_maps(const MapPipeline& input);
\end{lstlisting}
Procède par petites tâches :
\begin{enumerate}
\item Pour une entrée vide, retourne un pipeline vide.
\item Pars d'un morphisme scalaire accumulé égal à l'identité.
\item Pour chaque étape \code{stage}, mets à jour l'accumulateur par
\code{compose(stage, accumulated)}. Vérifie cet ordre à la main sur
\code{[f, g]} avant d'écrire la boucle.
\item Pour un pipeline non vide, retourne un pipeline contenant un seul
morphisme : celui que tu as accumulé. Dans cette première version, conserve
cette étape même si son chemin scalaire est vide.
\item Vérifie que l'entrée reste inchangée et qu'aucune liste de données
n'est évaluée pendant la fusion : la fonction manipule des descriptions.
\end{enumerate}

Garde d'abord la simplification scalaire de l'étape 5 séparée de cette
transformation. Quand chacune fonctionne, tu pourras simplifier le morphisme
fusionné. Tu sauras alors quelle transformation explique chaque changement.

\subsection{7.3 — Vérifier le résultat et l'ordre}
Construis \(f=A_1\), \(g=S\) et la liste \([0,1,2,3]\). Tu dois obtenir :
\begin{lstlisting}[style=plain]
Avant   : Map(Add(1)) >> Map(Square)
Apres   : Map(Add(1) >> Square)
Entree  : [0, 1, 2, 3]
Resultat: [1, 4, 9, 16]
\end{lstlisting}
Les résultats avant et après doivent être égaux. Avec l'ordre inverse
\code{[g, f]}, le résultat doit être \([1,2,5,10]\). Vérifie également
un pipeline vide, un pipeline d'une seule étape, trois étapes et une liste vide.
Teste expressément un pipeline contenant seulement \code{identity()} :
après fusion, il doit garder une étape, même s'il calcule la même fonction
qu'un pipeline vide. Leurs descriptions restent distinctes.

\newpage
\subsection{7.4 — Comprendre le gain sans brouiller la mesure}
Pour une liste de longueur \(N\), un pipeline de \(k>0\) étapes effectue
\(k\) parcours \code{map_values}, contre un après fusion. Pour notre
interpréteur simple, le bilan des \emph{étapes map}, hors copie initiale
commune de l'entrée, est le suivant :

\begin{center}
\begin{tabular}{lcc}
\toprule
Quantité & Avant & Après \\
\midrule
Parcours \code{map_values} & \(k\) & \(1\) \\
Visites d'éléments par ces parcours & \(kN\) & \(N\) \\
Listes produites par ces parcours & \(k\) & \(1\) \\
Listes intermédiaires issues des \code{map} & \(k-1\) & \(0\) \\
\bottomrule
\end{tabular}
\end{center}

Le nombre d'opérations scalaires reste le même si l'on ne simplifie pas
le morphisme composé : chaque élément subit toujours \(f\), puis \(g\), etc.
Les gains portent ici sur les parcours et les listes matérialisées.
Une liste vide ne contient aucun élément à visiter et peut n'exiger aucune
allocation. « Liste produite » n'est donc pas un comptage exact d'appels
à l'allocateur.

Commence par afficher les nombres d'étapes des descriptions. Si tu veux
vérifier les parcours réellement exécutés, ajoute ensuite un compteur dans
le pilote de l'interpréteur, autour des appels à \code{map_values}.
Ce compteur est un outil d'observation ; il ne fait pas partie du résultat
mathématique des morphismes.

\begin{vigilance}
Ne conclus pas encore « le programme est deux fois plus rapide ».
L'interpréteur a ses propres coûts, notamment la lecture des descriptions.
Un chronométrage exige un protocole distinct. La conclusion établie ici est
précise : mêmes valeurs, moins de parcours et de listes intermédiaires dans
cette exécution.
\end{vigilance}

\subsection{Questions à te poser}
Quelle différence y a-t-il entre \code{Morphism::steps} et
\code{MapPipeline::stages} ? Pourquoi la loi de foncteur autorise-t-elle
la fusion, mais pas l'inversion de \(f\) et \(g\) ? Pourquoi la pureté
est-elle essentielle si l'ordre des appels entre éléments change ?

\begin{validation}
\checkpoint La fusion conserve l'ordre des fonctions et les résultats.\par
\checkpoint Un pipeline non vide a exactement une étape après fusion.\par
\checkpoint Le pipeline vide reste vide et tous les cas limites passent.\par
\checkpoint Tu peux citer la loi qui justifie la transformation.\par
\checkpoint Tu peux expliquer le gain sans confondre parcours et opérations scalaires.
\end{validation}

\etape{8}{Consolider le C++ sans perdre le modèle}
\label{sec:consolider}
\begin{objectif}
Rendre ton petit langage plus lisible et plus difficile à utiliser
incorrectement, tout en conservant exactement les comportements déjà testés.
Cette étape prépare une extension du projet sans en changer le sens.
\end{objectif}

\textbf{C++ à travailler :} interfaces, invariants, compilation séparée,
encapsulation et responsabilité des types.

\subsection{À faire, dans cet ordre}
\begin{enumerate}
\item Sépare les déclarations, les définitions et le programme de démonstration.
Une organisation suffisante est :
\begin{lstlisting}[style=plain]
CatOpt/
  catopt.hpp     declarations et types du langage
  catopt.cpp     evaluation, composition et transformations
  tests.hpp      declaration de run_tests
  tests.cpp      definition de run_tests et assertions
  main.cpp       appel des tests et demonstration finale
  build/         fichiers produits par la compilation
\end{lstlisting}
Protège chaque en-tête contre les inclusions multiples, par exemple avec
\code{#pragma once}. Inclus \code{catopt.hpp} dans les fichiers qui utilisent
ses types et \code{tests.hpp} dans \code{main.cpp} et \code{tests.cpp}.
\item Compile les trois fichiers sources ensemble :
\begin{lstlisting}[style=plain]
clang++ -std=c++20 -Wall -Wextra -Wpedantic -g -O0 \
    main.cpp catopt.cpp tests.cpp -o build/catopt
\end{lstlisting}
Une déclaration dans l'en-tête permet d'utiliser une fonction ; sa définition
dans \code{catopt.cpp} fournit le code associé lors de l'édition de liens.
\item Regroupe les noms du langage dans \code{namespace catopt}.
Actualise les déclarations, les définitions et leurs utilisations.
\item Renforce \code{Primitive} : empêche désormais la construction de
combinaisons arbitraires de champs. Tu peux en faire une classe avec
constructeur privé, fabriques publiques statiques et accès en lecture.
Adapte les appels et les tests avant de modifier un autre type.
\item Pour \code{Morphism}, décide quelles opérations doivent avoir accès
au stockage. N'ajoute pas automatiquement un setter pour chaque champ.
Une interface de lecture et les opérations de construction peuvent suffire.
\end{enumerate}

\begin{theorie}
L'encapsulation protège les représentations sur lesquelles tes lois reposent.
Par exemple, une primitive bien construite désigne toujours une opération
reconnue, et un morphisme garde un ordre d'exécution explicite.
Une classe n'est pas à elle seule une preuve mathématique ; elle aide à
faire respecter les hypothèses de ton modèle par le code appelant.
\end{theorie}

\subsection{Une lecture mémoire utile}
Choisis une fonction et explique la durée de vie de chaque objet qu'elle
utilise. Identifie une copie nécessaire et une copie éventuellement évitable.
Ne remplace pas chaque retour par \code{std::move} : le retour par valeur
permet déjà, selon la situation, l'élision de copie ou le déplacement.
Comprends d'abord le contrat de propriété.

Avec Clang ou GCC sur une plateforme qui les prend en charge, tu peux aussi
recompiler avec \code{-fsanitize=address,undefined -fno-omit-frame-pointer}
et exécuter tes exemples. Ces outils recherchent certaines erreurs mémoire
et certains comportements indéfinis ; ils ne prouvent pas les lois de CatOpt.

\begin{validation}
\checkpoint Le programme compile depuis ses différents fichiers.\par
\checkpoint Tous les exemples précédents ont le même résultat.\par
\checkpoint Tu peux expliquer les conventions devenues des invariants protégés.\par
\checkpoint Tu sais qui possède chaque vecteur et pourquoi ses éléments restent valides.
\end{validation}

\clearpage
\section{La démonstration qui termine ce premier volume}
\label{sec:demo}
Prépare un exécutable qui raconte la transformation en quatre temps.
Son affichage peut ressembler à ceci :

\begin{lstlisting}[style=plain]
1. Un morphisme scalaire
   f = Add(1)
   g = Square
   g o f = Add(1) >> Square
   (g o f)(3) = 16

2. Une reecriture arithmetique
   Avant : Add(1) >> Add(2) >> Square
   Apres : Add(3) >> Square
   Loi   : A_b o A_a = A_(a+b)
   Entree 3 : 36 dans les deux cas

3. Une reecriture par fonctorialite
   Avant : Map(Add(1)) >> Map(Square)
   Apres : Map(Add(1) >> Square)
   Loi   : F(g) o F(f) = F(g o f)
   Entree : [0, 1, 2, 3]
   Resultat : [1, 4, 9, 16]
   Parcours map : 2 -> 1

4. Verification
   Tous les tests executes sont passes.
\end{lstlisting}

Cet affichage est une démonstration explicative. Ce n'est pas encore un
certificat vérifié indépendamment : les règles et l'interpréteur restent
des éléments de confiance de ton programme.

\subsection{Les cinq questions auxquelles tu dois savoir répondre}
\begin{enumerate}
\item Quel est l'objet de la catégorie scalaire, et que représentent ses morphismes ?
\item Pourquoi l'identité est-elle une suite vide ?
\item Pourquoi \code{compose(g, f)} stocke-t-il d'abord \code{f} ?
\item Quelle différence sépare une loi de catégorie, une loi arithmétique
et une loi de foncteur ?
\item Qu'est-ce qui est préservé par l'optimiseur, et quel coût diminue
effectivement dans ta première démonstration ?
\end{enumerate}

\begin{objectif}
À ce stade, CatOpt possède une syntaxe, une sémantique, une composition,
des tests de lois et deux sortes de réécritures. Tu as travaillé les bases
du C++ dans un programme cohérent, et tu peux expliquer chaque transformation
avec les objets mathématiques qui l'ont motivée.
\end{objectif}

\clearpage
\section{La suite : choisir une nouvelle difficulté à la fois}
\label{sec:suite}
Le noyau par chemins est volontairement limité : il décrit des calculs
linéaires sur un seul type scalaire, puis des parcours de listes. Tu peux
le faire évoluer lorsque tu sais expliquer ses limites concrètes.

\begin{center}
\small
\begin{tabularx}{\linewidth}{@{}>{\raggedright\arraybackslash}p{0.24\linewidth}
  >{\raggedright\arraybackslash}X >{\raggedright\arraybackslash}X@{}}
\toprule
Prochain besoin & Notions C++ & Question catégorique \\
\midrule
Plusieurs objets et types intermédiaires & Types explicites, validation,
\code{std::variant} & Quand \(g\circ f\) est-il bien typé ? \\
\addlinespace
Une syntaxe avec \code{Compose}, \code{Map} et produits imbriqués & Arbres,
visiteurs, récursivité, \code{std::unique_ptr} & Comment conserver la
structure d'un morphisme composé ? \\
\addlinespace
Des sous-expressions partagées & Identifiants, arène, propriété partagée
si nécessaire & Comment plusieurs chemins réutilisent-ils une description ? \\
\addlinespace
Une interface typée à la compilation & Templates, déduction, concepts &
Peut-on refuser une composition dont les objets intermédiaires diffèrent ? \\
\addlinespace
Deux calculs indépendants & Types produits, représentation de graphes &
Que signifie \(f\otimes g\), et quand la loi d'interchange est-elle valide ? \\
\addlinespace
Explorer plusieurs réécritures & Hachage, union-find, e-graphs & Comment
représenter plusieurs morphismes égaux modulo des équations ? \\
\bottomrule
\end{tabularx}
\end{center}

\subsection{Une première expérience de templates, au bon moment}
Quand le noyau fonctionne, écris à part une fonction de composition de deux
objets appelables, paramétrée par leurs types. Compare-la à
\code{compose(const Morphism&, const Morphism&)}.
La première construit un calcul exécutable ; la seconde construit une
description inspectable. Comprendre cette différence t'évitera de perdre
la capacité de réécriture en remplaçant trop tôt toutes les descriptions
par des fonctions opaques.

\subsection{Un premier backend}
Une autre suite naturelle consiste à produire un fichier C++ à partir
du morphisme fusionné. Tu pourras alors vérifier qu'une seule boucle est
générée, compiler le fichier et comparer ses résultats à ceux de
l'interpréteur de référence. Le passage du modèle mathématique au code
généré deviendra une nouvelle obligation de correction.

\textbf{Conseil de progression :} choisis une seule de ces extensions pour
le prochain jalon. Le but reste de comprendre du C++ à travers les
constructions catégoriques que tu as envie d'explorer.

\clearpage
\section{Repères de travail et ressources}
\label{sec:reperes}
\subsection{Petit dictionnaire du projet}
\begin{description}[style=nextline,leftmargin=0pt,labelindent=0pt,labelsep=0pt,
  labelwidth=0pt,font=\bfseries]
\item[Objet] Un domaine de valeurs. Ici, \(U\), puis \(\List(U)\).
\item[Morphisme] Une transformation avec un domaine et un codomaine.
Nos morphismes scalaires vont tous de \(U\) vers \(U\).
\item[Générateur] Une opération élémentaire, ici \(A_c\) ou \(S\).
\item[Composition] Enchaînement : \(g\circ f\) exécute \(f\), puis \(g\).
\item[Identité] Morphisme qui laisse la valeur inchangée ; chemin vide
dans notre syntaxe scalaire.
\item[Interprétation] Passage d'une description \(p\) à la fonction
\(\sem{p}\) qu'elle représente.
\item[Foncteur] Application sur les objets et les morphismes qui préserve
l'identité et la composition.
\item[Réécriture] Remplacement d'une description par une autre selon une
règle justifiée, éventuellement à l'intérieur d'un programme.
\item[Modèle de coût] Critère explicite pour comparer les descriptions :
nombre de primitives, de parcours, etc.
\end{description}

\subsection{La fiche à remplir après chaque séance}
\begin{enumerate}
\item \textbf{Dans CatOpt :} qu'est-ce que mon programme sait faire de plus ?
\item \textbf{En catégories :} quel objet, morphisme ou quelle loi ai-je manipulé ?
\item \textbf{En C++ :} quelle notion ai-je utilisée, et pourquoi ?
\item \textbf{Vérification :} quel exemple pouvait révéler une erreur ?
\item \textbf{Prochaine petite tâche :} quelle fonction ou quel cas vais-je traiter ?
\end{enumerate}

\subsection{Lectures ciblées, facultatives}
\begin{itemize}
\item Pour vérifier la sémantique des entiers :
\href{https://eel.is/c++draft/basic.fundamental}{brouillon de travail du standard C++, types fondamentaux},
notamment le paragraphe sur l'arithmétique non signée. Ce document technique
évolue ; le comportement modulaire utilisé ici est établi en C++20.
\item Pour voir une bibliothèque qui manipule objets, morphismes et
diagrammes : \href{https://docs.discopy.org/en/main/index.html}{documentation de DisCoPy}.
Elle propose notamment des diagrammes, des compositions et des interprétations
par foncteurs. Elle est une source d'exemples, pas un prérequis pour ce guide.
\item Après les premières réécritures, pour découvrir la saturation d'égalités :
\href{https://www.mwillsey.com/papers/egg}{Willsey et al., \emph{egg: Fast and Extensible Equality Saturation}, POPL 2021}.
Cette lecture prépare une extension ultérieure.
\end{itemize}

\subsection{Compiler ce document}
Le fichier \code{CatOpt-guide-demarrage.tex} est autonome et peut être
importé dans Overleaf avec le compilateur pdfLaTeX. Avec une installation
TeX Live ou MacTeX incluant les paquets utilisés, place-toi dans son dossier :
\begin{lstlisting}[style=plain]
latexmk -pdf -interaction=nonstopmode -halt-on-error \
    CatOpt-guide-demarrage.tex
\end{lstlisting}

\textbf{Ta première séance :} termine l'étape 0, puis écris les trois
fonctions scalaires de l'étape 1. Arrête-toi quand tu sais expliquer pourquoi
\(S\circ A_1\) envoie \(3\) sur \(16\), et pourquoi l'ordre inverse donne \(10\).

\end{document}
